\begin{table*}[t]\footnotesize\centering\renewcommand{\arraystretch}{0.8}
\caption{Main results: Spearman and top-100 recall (mean $\pm$ std over \rhval{const/n_seeds} draws) on all non-training variants (rand) or on Hamming-distance$\ge3$ variants (dbl). Bold: best mean in the task; underline: second.}\label{tab:main}
\begin{tabular}{llccc}
\toprule
Method & Task & spearman $\uparrow$ & top100\_recall $\uparrow$ & $n$ \\
\midrule
One-hot ridge (reimplemented) & rand\_48 & 0.182 $\pm$ 0.052 & \textbf{0.006 $\pm$ 0.009} & 5 \\
GP Hamming kernel (reimplemented) & rand\_48 & \textbf{0.202 $\pm$ 0.051} & 0.004 $\pm$ 0.009 & 5 \\
CNN (reimplemented) & rand\_48 & 0.166 $\pm$ 0.072 & 0.004 $\pm$ 0.009 & 5 \\
\textbf{One-hot + pairwise ridge} & rand\_48 & \underline{0.188 $\pm$ 0.043} & \underline{0.006 $\pm$ 0.009} & 5 \\
\midrule
One-hot ridge (reimplemented) & rand\_96 & 0.231 $\pm$ 0.036 & \underline{0.016 $\pm$ 0.018} & 5 \\
GP Hamming kernel (reimplemented) & rand\_96 & \underline{0.236 $\pm$ 0.030} & 0.010 $\pm$ 0.017 & 5 \\
CNN (reimplemented) & rand\_96 & \textbf{0.287 $\pm$ 0.042} & \textbf{0.022 $\pm$ 0.016} & 5 \\
\textbf{One-hot + pairwise ridge} & rand\_96 & 0.229 $\pm$ 0.035 & 0.012 $\pm$ 0.016 & 5 \\
\midrule
One-hot ridge (reimplemented) & rand\_384 & 0.381 $\pm$ 0.037 & 0.024 $\pm$ 0.027 & 5 \\
GP Hamming kernel (reimplemented) & rand\_384 & 0.389 $\pm$ 0.042 & \textbf{0.036 $\pm$ 0.025} & 5 \\
CNN (reimplemented) & rand\_384 & \textbf{0.444 $\pm$ 0.023} & 0.030 $\pm$ 0.014 & 5 \\
\textbf{One-hot + pairwise ridge} & rand\_384 & \underline{0.396 $\pm$ 0.052} & \underline{0.034 $\pm$ 0.028} & 5 \\
\midrule
One-hot ridge (reimplemented) & rand\_2000 & \underline{0.480 $\pm$ 0.017} & 0.084 $\pm$ 0.080 & 5 \\
GP Hamming kernel (reimplemented) & rand\_2000 & 0.451 $\pm$ 0.006 & \underline{0.152 $\pm$ 0.062} & 5 \\
CNN (reimplemented) & rand\_2000 & \textbf{0.506 $\pm$ 0.015} & \textbf{0.216 $\pm$ 0.109} & 5 \\
\textbf{One-hot + pairwise ridge} & rand\_2000 & 0.478 $\pm$ 0.009 & 0.152 $\pm$ 0.063 & 5 \\
\midrule
One-hot ridge (reimplemented) & dbl\_48 & \textbf{0.212 $\pm$ 0.063} & \textbf{0.018 $\pm$ 0.027} & 5 \\
GP Hamming kernel (reimplemented) & dbl\_48 & \underline{0.199 $\pm$ 0.057} & 0.014 $\pm$ 0.022 & 5 \\
CNN (reimplemented) & dbl\_48 & 0.186 $\pm$ 0.081 & 0.012 $\pm$ 0.027 & 5 \\
\textbf{One-hot + pairwise ridge} & dbl\_48 & 0.192 $\pm$ 0.048 & \underline{0.016 $\pm$ 0.023} & 5 \\
\midrule
One-hot ridge (reimplemented) & dbl\_96 & 0.262 $\pm$ 0.038 & 0.006 $\pm$ 0.009 & 5 \\
GP Hamming kernel (reimplemented) & dbl\_96 & 0.252 $\pm$ 0.046 & \textbf{0.012 $\pm$ 0.011} & 5 \\
CNN (reimplemented) & dbl\_96 & \textbf{0.289 $\pm$ 0.046} & 0.008 $\pm$ 0.008 & 5 \\
\textbf{One-hot + pairwise ridge} & dbl\_96 & \underline{0.263 $\pm$ 0.045} & \underline{0.010 $\pm$ 0.010} & 5 \\
\midrule
One-hot ridge (reimplemented) & dbl\_384 & \underline{0.320 $\pm$ 0.046} & \underline{0.020 $\pm$ 0.016} & 5 \\
GP Hamming kernel (reimplemented) & dbl\_384 & 0.299 $\pm$ 0.040 & 0.014 $\pm$ 0.013 & 5 \\
CNN (reimplemented) & dbl\_384 & \textbf{0.332 $\pm$ 0.053} & \textbf{0.028 $\pm$ 0.013} & 5 \\
\textbf{One-hot + pairwise ridge} & dbl\_384 & 0.247 $\pm$ 0.054 & 0.014 $\pm$ 0.013 & 5 \\
\midrule
One-hot ridge (reimplemented) & dbl\_2000 & \underline{0.354 $\pm$ 0.005} & 0.062 $\pm$ 0.011 & 5 \\
GP Hamming kernel (reimplemented) & dbl\_2000 & 0.322 $\pm$ 0.013 & 0.114 $\pm$ 0.009 & 5 \\
CNN (reimplemented) & dbl\_2000 & \textbf{0.387 $\pm$ 0.034} & \textbf{0.208 $\pm$ 0.023} & 5 \\
\textbf{One-hot + pairwise ridge} & dbl\_2000 & 0.196 $\pm$ 0.010 & \underline{0.128 $\pm$ 0.016} & 5 \\
\bottomrule
\end{tabular}

\end{table*}

\begin{figure}[t]\centering
\includegraphics[width=\columnwidth]{learning_curves.pdf}
\caption{Learning curves (mean $\pm$ std over draws) for the four systems under random (left) and single/double-mutant (right) training. Top: Spearman; bottom: top-100 recall.}\label{fig:lc}
\end{figure}

Table~\ref{tab:main} and Fig.~\ref{fig:lc} give the main results. Differences are judged against the seed-to-seed standard deviation, which is about \rhval{main/one-hot-ridge-reimplemented/rand_48/spearman/std:3} Spearman at $N{=}48$ and shrinks to \rhval{main/one-hot-ridge-reimplemented/rand_2000/spearman/std:3} at $N{=}2000$. Quoted $p$ are paired $t$-test values over five draws from \texttt{rh compare}, uncorrected.

\subsection{H1: pairwise epistasis features}
\emph{Prediction.} Pairwise features help only at $N\ge384$. \emph{Observation.} With random training, pairwise and additive ridge are indistinguishable at every size: Spearman \rhval{main/one-hot-+-pairwise-ridge/rand_48/spearman/mean:3} versus \rhval{main/one-hot-ridge-reimplemented/rand_48/spearman/mean:3} ($N{=}48$, $p=\rhval{cmp/main/one-hot-ridge-reimplemented/rand_48/spearman/paired_p}$), \rhval{main/one-hot-+-pairwise-ridge/rand_96/spearman/mean:3} versus \rhval{main/one-hot-ridge-reimplemented/rand_96/spearman/mean:3} ($N{=}96$), \rhval{main/one-hot-+-pairwise-ridge/rand_384/spearman/mean:3} versus \rhval{main/one-hot-ridge-reimplemented/rand_384/spearman/mean:3} ($N{=}384$, $p=\rhval{cmp/main/one-hot-ridge-reimplemented/rand_384/spearman/paired_p}$) and \rhval{main/one-hot-+-pairwise-ridge/rand_2000/spearman/mean:3} versus \rhval{main/one-hot-ridge-reimplemented/rand_2000/spearman/mean:3} ($N{=}2000$, $p=\rhval{cmp/main/one-hot-ridge-reimplemented/rand_2000/spearman/paired_p}$). \emph{Verdict: refuted} by the pre-stated rule (the pairwise mean is above ridge at $N{=}48$ and not above it at $N{=}2000$), but because every difference is far inside the noise this is better read as ``no detectable effect of pairwise features with random training'' than as a reversed effect. Top-100 recall at $N{=}2000$ is higher for pairwise ridge than for additive ridge (\rhval{main/one-hot-+-pairwise-ridge/rand_2000/top100_recall/mean:3} versus \rhval{main/one-hot-ridge-reimplemented/rand_2000/top100_recall/mean:3}, paired $p=\rhval{cmp/main/one-hot-ridge-reimplemented/rand_2000/top100_recall/paired_p}$), the only hint of a benefit; we did not register this as a test, and the log-target ablation (Table~\ref{tab:logt}) shows that the ordering depends on the target transform.
\emph{Failure case.} Under single/double-mutant training pairwise ridge is the worst system at the two largest sizes: at $N{=}2000$ it reaches \rhval{main/one-hot-+-pairwise-ridge/dbl_2000/spearman/mean:3} against \rhval{main/one-hot-ridge-reimplemented/dbl_2000/spearman/mean:3} for the additive model ($p=\rhval{cmp/main/one-hot-ridge-reimplemented/dbl_2000/spearman/paired_p}$), and at $N{=}384$ \rhval{main/one-hot-+-pairwise-ridge/dbl_384/spearman/mean:3} against \rhval{main/one-hot-ridge-reimplemented/dbl_384/spearman/mean:3} ($p=\rhval{cmp/main/one-hot-ridge-reimplemented/dbl_384/spearman/paired_p}$). The model fits the training doubles but its learned interactions do not transfer to variants three or four mutations away (Table~\ref{tab:hd}); we have not isolated the cause, but a plausible one is that, with only wild-type, single and double variants, the pair indicators of site pairs that are both wild type co-vary with mutations at the other two sites and absorb effects that do not generalise.

\subsection{H2: Gaussian process at small N}
\emph{Prediction.} GP $\ge$ ridge at $N\le96$. \emph{Observation.} GP means are higher at $N{=}48$ (\rhval{main/gp-hamming-kernel-reimplemented/rand_48/spearman/mean:3} versus \rhval{main/one-hot-ridge-reimplemented/rand_48/spearman/mean:3}) and $N{=}96$ (\rhval{main/gp-hamming-kernel-reimplemented/rand_96/spearman/mean:3} versus \rhval{main/one-hot-ridge-reimplemented/rand_96/spearman/mean:3}). \emph{Verdict: supported in the mean, inconclusive statistically.} The GP is not a better model at larger N: at $N{=}2000$ it is the weakest system (\rhval{main/gp-hamming-kernel-reimplemented/rand_2000/spearman/mean:3}) and clearly below pairwise ridge ($p=\rhval{cmp/main/gp-hamming-kernel-reimplemented/rand_2000/spearman/paired_p}$).

\subsection{H3: small CNN}
\emph{Prediction.} The CNN is worse than the best linear or kernel model at $N\le384$ and does not clearly exceed it at $N{=}2000$. \emph{Observation.} The CNN has the lowest mean at $N{=}48$ (\rhval{main/cnn-reimplemented/rand_48/spearman/mean:3}) but the highest at $N{=}96$ (\rhval{main/cnn-reimplemented/rand_96/spearman/mean:3}; against pairwise ridge $p=\rhval{cmp/main/cnn-reimplemented/rand_96/spearman/paired_p}$), at $N{=}384$ (\rhval{main/cnn-reimplemented/rand_384/spearman/mean:3}; against pairwise ridge $p=\rhval{cmp/main/cnn-reimplemented/rand_384/spearman/paired_p}$) and at $N{=}2000$ (\rhval{main/cnn-reimplemented/rand_2000/spearman/mean:3}; against pairwise ridge $p=\rhval{cmp/main/cnn-reimplemented/rand_2000/spearman/paired_p}$), and at $N{=}2000$ it has the highest top-100 recall (\rhval{main/cnn-reimplemented/rand_2000/top100_recall/mean:3}, with a large spread, std \rhval{main/cnn-reimplemented/rand_2000/top100_recall/std:3}). The same ordering holds under single/double-mutant training at $N{=}2000$ (\rhval{main/cnn-reimplemented/dbl_2000/spearman/mean:3}). \emph{Verdict: refuted} as stated. Two caveats limit what this says about CNNs: the model was tuned more than the others, and Sec.~\ref{sec:ablations} shows that the advantage mostly vanishes when all models regress a log-transformed fitness.

\subsection{H4: random versus single/double-mutant training}
Table~\ref{tab:gap} gives the difference of mean Spearman on test variants of Hamming distance 3 and 4 between random and single/double training at equal $N$ (positive: random training better). For $N\ge384$ random training is better or equal for every system on both distances in the mean, but the prediction fails at small N: at $N\le96$ the Hamming-distance-3 difference is negative for every system (for example \rhval{derived/derived-analysis/rand_384/gap_hd3_ridge_48/mean:3} for additive ridge at $N{=}48$), i.e.\ single/double training generalises better to distance-3 variants than a random draw of the same size. A possible reason, which we did not test, is that a small random draw consists almost entirely of distance-4 variants and so contains few informative neighbours of the test variants. \emph{Verdict: refuted} as stated (``for every system'' fails at $N\le96$), supported for $N\ge384$.

\subsection{H5: top-100 recall}
The largest mean top-100 recall at any $N\le384$ is \rhval{main/gp-hamming-kernel-reimplemented/rand_384/top100_recall/mean:3} (GP, rand$_{384}$), so no model recovers more than a few of the 100 best variants from a few hundred labels. \emph{Verdict: supported.} At $N{=}2000$ recall rises to \rhval{main/cnn-reimplemented/rand_2000/top100_recall/mean:3} for the CNN and \rhval{main/one-hot-+-pairwise-ridge/rand_2000/top100_recall/mean:3} for pairwise ridge; the differences between systems are large relative to their means but the std across draws is of the same order (e.g.\ \rhval{main/one-hot-+-pairwise-ridge/rand_2000/top100_recall/std:3} for pairwise ridge), so only the gap between the CNN or kernel/pairwise models and additive ridge (\rhval{main/one-hot-ridge-reimplemented/rand_2000/top100_recall/mean:3}) is suggestive.

\begin{table*}[t]\footnotesize\centering\renewcommand{\arraystretch}{0.8}
\caption{Spearman by Hamming distance of the test variant (mean $\pm$ std; HD3, HD4), main group. Test variants at Hamming distance 2 are omitted because fewer than the minimum needed exist in the dbl tasks.}\label{tab:hd}
\resizebox{0.6\textwidth}{!}{\begin{tabular}{llccc}
\toprule
Method & Task & spearman\_hd3 $\uparrow$ & spearman\_hd4 $\uparrow$ & $n$ \\
\midrule
One-hot ridge (reimplemented) & rand\_48 & 0.257 $\pm$ 0.181 & 0.145 $\pm$ 0.064 & 5 \\
GP Hamming kernel (reimplemented) & rand\_48 & \textbf{0.284 $\pm$ 0.131} & \textbf{0.161 $\pm$ 0.089} & 5 \\
CNN (reimplemented) & rand\_48 & 0.169 $\pm$ 0.174 & \underline{0.152 $\pm$ 0.086} & 5 \\
\textbf{One-hot + pairwise ridge} & rand\_48 & \underline{0.265 $\pm$ 0.152} & 0.151 $\pm$ 0.075 & 5 \\
\midrule
One-hot ridge (reimplemented) & rand\_96 & 0.294 $\pm$ 0.060 & 0.209 $\pm$ 0.041 & 5 \\
GP Hamming kernel (reimplemented) & rand\_96 & \underline{0.316 $\pm$ 0.040} & \underline{0.209 $\pm$ 0.042} & 5 \\
CNN (reimplemented) & rand\_96 & \textbf{0.340 $\pm$ 0.039} & \textbf{0.268 $\pm$ 0.047} & 5 \\
\textbf{One-hot + pairwise ridge} & rand\_96 & 0.296 $\pm$ 0.067 & 0.205 $\pm$ 0.037 & 5 \\
\midrule
One-hot ridge (reimplemented) & rand\_384 & 0.535 $\pm$ 0.026 & 0.322 $\pm$ 0.039 & 5 \\
GP Hamming kernel (reimplemented) & rand\_384 & 0.538 $\pm$ 0.035 & 0.331 $\pm$ 0.044 & 5 \\
CNN (reimplemented) & rand\_384 & \textbf{0.570 $\pm$ 0.036} & \textbf{0.395 $\pm$ 0.032} & 5 \\
\textbf{One-hot + pairwise ridge} & rand\_384 & \underline{0.549 $\pm$ 0.034} & \underline{0.338 $\pm$ 0.056} & 5 \\
\midrule
One-hot ridge (reimplemented) & rand\_2000 & 0.633 $\pm$ 0.008 & \underline{0.421 $\pm$ 0.021} & 5 \\
GP Hamming kernel (reimplemented) & rand\_2000 & 0.616 $\pm$ 0.016 & 0.391 $\pm$ 0.010 & 5 \\
CNN (reimplemented) & rand\_2000 & \textbf{0.649 $\pm$ 0.015} & \textbf{0.455 $\pm$ 0.018} & 5 \\
\textbf{One-hot + pairwise ridge} & rand\_2000 & \underline{0.643 $\pm$ 0.010} & 0.419 $\pm$ 0.013 & 5 \\
\midrule
One-hot ridge (reimplemented) & dbl\_48 & \textbf{0.438 $\pm$ 0.030} & \underline{0.149 $\pm$ 0.076} & 5 \\
GP Hamming kernel (reimplemented) & dbl\_48 & \underline{0.413 $\pm$ 0.046} & \textbf{0.151 $\pm$ 0.064} & 5 \\
CNN (reimplemented) & dbl\_48 & 0.320 $\pm$ 0.083 & 0.147 $\pm$ 0.086 & 5 \\
\textbf{One-hot + pairwise ridge} & dbl\_48 & 0.376 $\pm$ 0.068 & 0.142 $\pm$ 0.057 & 5 \\
\midrule
One-hot ridge (reimplemented) & dbl\_96 & \underline{0.469 $\pm$ 0.030} & 0.209 $\pm$ 0.041 & 5 \\
GP Hamming kernel (reimplemented) & dbl\_96 & 0.459 $\pm$ 0.032 & 0.213 $\pm$ 0.052 & 5 \\
CNN (reimplemented) & dbl\_96 & \textbf{0.480 $\pm$ 0.033} & \textbf{0.250 $\pm$ 0.048} & 5 \\
\textbf{One-hot + pairwise ridge} & dbl\_96 & 0.455 $\pm$ 0.048 & \underline{0.217 $\pm$ 0.045} & 5 \\
\midrule
One-hot ridge (reimplemented) & dbl\_384 & \underline{0.525 $\pm$ 0.023} & \underline{0.262 $\pm$ 0.050} & 5 \\
GP Hamming kernel (reimplemented) & dbl\_384 & 0.501 $\pm$ 0.027 & 0.261 $\pm$ 0.042 & 5 \\
CNN (reimplemented) & dbl\_384 & \textbf{0.539 $\pm$ 0.043} & \textbf{0.290 $\pm$ 0.058} & 5 \\
\textbf{One-hot + pairwise ridge} & dbl\_384 & 0.324 $\pm$ 0.053 & 0.231 $\pm$ 0.057 & 5 \\
\midrule
One-hot ridge (reimplemented) & dbl\_2000 & \underline{0.549 $\pm$ 0.002} & 0.299 $\pm$ 0.005 & 5 \\
GP Hamming kernel (reimplemented) & dbl\_2000 & 0.446 $\pm$ 0.007 & \underline{0.310 $\pm$ 0.014} & 5 \\
CNN (reimplemented) & dbl\_2000 & \textbf{0.591 $\pm$ 0.033} & \textbf{0.347 $\pm$ 0.030} & 5 \\
\textbf{One-hot + pairwise ridge} & dbl\_2000 & 0.245 $\pm$ 0.005 & 0.189 $\pm$ 0.013 & 5 \\
\bottomrule
\end{tabular}
}
\end{table*}

\begin{table}[t]\footnotesize\centering\renewcommand{\arraystretch}{0.8}
\caption{Mean Spearman on Hamming-distance-3 and 4 test variants: random minus single/double training at equal $N$.}\label{tab:gap}
\resizebox{\columnwidth}{!}{\begin{tabular}{llrrrr}
\toprule
System & Test HD & N=48 & N=96 & N=384 & N=2000 \\
\midrule
One-hot ridge & 3 & -0.181 & -0.175 & +0.010 & +0.084 \\
One-hot ridge & 4 & -0.004 & +0.000 & +0.059 & +0.122 \\
Pairwise ridge & 3 & -0.111 & -0.159 & +0.225 & +0.398 \\
Pairwise ridge & 4 & +0.009 & -0.012 & +0.106 & +0.230 \\
GP & 3 & -0.129 & -0.143 & +0.037 & +0.170 \\
GP & 4 & +0.010 & -0.004 & +0.069 & +0.081 \\
CNN & 3 & -0.151 & -0.140 & +0.031 & +0.058 \\
CNN & 4 & +0.005 & +0.018 & +0.105 & +0.107 \\
\bottomrule
\end{tabular}
}
\end{table}

\subsection{Where accuracy is lost}
Table~\ref{tab:hd} shows that all models are worst on distance-4 variants, which are the \rhval{const/n_hd4} of \rhval{const/n_variants} variants furthest from wild type: for example the CNN at rand$_{2000}$ reaches \rhval{main/cnn-reimplemented/rand_2000/spearman_hd3/mean:3} on distance-3 variants but \rhval{main/cnn-reimplemented/rand_2000/spearman_hd4/mean:3} on distance-4 variants. The overall Spearman is therefore dominated by the distance-4 majority of the landscape.
